\section{实数系一些概念和问题}

\subsection*{实数的十进制表示}

对任意 $x \in \mathbb{R}$，$x \ge 0$，令 $m=[x]$. 原图记
\[
a_k=[10^k x]-[10^{k-1}x],
\]
并称其为“第 $k$ 位小数”. 此时
\[
x=\sup_{k \ge 1}\left(m+\sum_{i=1}^k \frac{a_i}{10^i}\right),
\]
记作
\[
x=m.a_1a_2\cdots a_k\cdots,
\qquad
-x=-m.a_1a_2\cdots a_k\cdots.
\]

\begin{notebox}
这里 $a_k$ 的定义按原图字面照录；但从通常十进制展开的写法看，式中疑似缺少一个系数 $10$. 由于原图未进一步说明，这里只做标注，不擅自改写. 
\end{notebox}

\subsection*{实数的 $p$ 进制表示}

原图进一步写到：对任意 $x \in \mathbb{R}$，$x \ge 0$，

当 $x \le 1$ 时，$N:=0$；当 $x>1$ 时，取 $N \in \mathbb{N}^*$ 使
\[
p^{N-1}<x<p^N.
\]

随后定义
\[
a_k=[p^{k-N}x]-[p^{k-1-N}x],
\]
并写出
\[
x=\sup_{n \to \infty}\sum_{k=-n}^{N} a_k p^k.
\]

记作
\[
x=a_Na_{N-1}\cdots a_0.a_{-1}a_{-2}\cdots \quad (\text{底 } p).
\]

\begin{notebox}
这一段中量词、下标范围和系数位置都较紧，已按可辨内容整理；其中 $a_k$ 的定义同样疑似缺少一个系数 $p$. 
\end{notebox}

\subsection*{定义 $n$ 次方根}

\begin{theorem}[n次方根的存在唯一性]
设 $a > 0$，$n \in \mathbb{N}^*$，则方程 $x^n = a$ 有唯一正解. 
\end{theorem}

\begin{proof}
$1^\circ$ （证明唯一性）

若 $s>t>0$，则 $s^n>t^n$. 因此若 $x_0^n=a$，则
\[
\begin{cases}
x'>x_0 \implies x'^n>x_0^n=a,\\
x'<x_0 \implies x'^n<x_0^n=a.
\end{cases}
\]
故 $x_0$ 唯一. 

$2^\circ$ （证明存在性）

构造
\[
E_a:=\{y \mid y>0,\ y^n\le a\}.
\]

(i) 证明 $E_a$ 非空且有上界. 

取 $s<1$ 且 $s<a$，则 $s^n<s<a$，故 $s \in E_a$，于是 $E_a \ne \varnothing$. 

若 $y \ge 1+a > 1$，则 $y^n \ge (1+a)^n > 1+a > a$，故 $y \notin E_a$. 
所以 $1+a$ 是 $E_a$ 的一个上界. 令
\[
x=\sup E_a.
\]

(ii) 证明 $x^n=a$. 

若 $x^n>a$，令
\[
\varepsilon=\frac{x^n-a}{(x+1)^n}\in(0,1),
\]
则
\[
\begin{aligned}
(x-\varepsilon)^n
&=x^n+\sum_{k=1}^n \binom{n}{k}x^{n-k}(-\varepsilon)^k \\
&>x^n-\sum_{k=1}^n \binom{n}{k}x^{n-k}\varepsilon^k \\
&>x^n-\varepsilon \sum_{k=1}^n \binom{n}{k}x^{n-k} \\
&=x^n-\varepsilon\big((x+1)^n-x^n\big) \\
&>x^n-\varepsilon(x+1)^n \\
&=x^n-(x^n-a)=a.
\end{aligned}
\]
于是 $x-\varepsilon$ 也是 $E_a$ 的上界，这与 $x$ 为上确界矛盾. 

若 $x^n<a$，令
\[
\varepsilon=
\begin{cases}
\dfrac{a-x^n}{(x+1)^n}, & \dfrac{a-x^n}{(x+1)^n}<1,\\[0.4em]
1, & \text{others},
\end{cases}
\qquad \varepsilon \in (0,1].
\]
则
\[
\begin{aligned}
(x+\varepsilon)^n
&=x^n+\sum_{k=1}^n \binom{n}{k}x^{n-k}\varepsilon^k \\
&\le x^n+\varepsilon \sum_{k=1}^n \binom{n}{k}x^{n-k} \\
&<x^n+\varepsilon\big((x+1)^n-x^n\big) \\
&<x^n+\varepsilon(x+1)^n \\
&\le x^n+(a-x^n)=a.
\end{aligned}
\]
故 $x+\varepsilon \in E_a$，与 $x$ 为上界矛盾. 

综上，$x^n=a$，故 $x$ 为方程 $x^n=a$ 的唯一正解. 
\end{proof}

定义 $a^{1/n}$ 为 $a$ 的 $n$ 次方根. 

\begin{theorem}[均值不等式]
设 $n \ge 2$，$a_1,a_2,\dots,a_n \ge 0$，则
\[
\sqrt[n]{\prod_{k=1}^n a_k}\le \frac{1}{n}\sum_{k=1}^n a_k.
\]
\end{theorem}

\begin{proof}
用数学归纳法证明. 

$1^\circ$ 当 $n=2$ 时，
\[
(\sqrt{a_1}-\sqrt{a_2})^2 \ge 0
\implies a_1+a_2 \ge 2\sqrt{a_1a_2}.
\]

$2^\circ$ 先证当 $n=2^k$ 时成立. 

设命题对 $n=2^k$ 成立，记 $m=2^k$，则
\[
\sqrt[m]{\prod_{i=1}^m a_i}\le \frac{1}{m}\sum_{i=1}^m a_i.
\]
当 $n=2^{k+1}=2m$ 时，
\[
\begin{aligned}
\sqrt[2m]{\prod_{i=1}^{2m} a_i}
&=\sqrt[m]{\prod_{i=1}^m \sqrt{a_i a_{i+m}}} \\
&\le \frac{1}{m}\sum_{i=1}^m \sqrt{a_i a_{i+m}} \\
&\le \frac{1}{m}\sum_{i=1}^m \frac{a_i+a_{i+m}}{2} \\
&=\frac{1}{2m}\sum_{i=1}^{2m} a_i.
\end{aligned}
\]
故对一切 $n=2^k$ 命题成立. 

$3^\circ$ 再证对一切 $n \ge 2$ 成立. 

设 $2^k \le n < 2^{k+1}$，记 $m=2^k$，则 $m \le n < 2m$. 令
\[
A=\frac{1}{n}\sum_{i=1}^n a_i.
\]
由 $2m$ 个数的情形可得
\[
\sqrt[2m]{\left(\prod_{i=1}^n a_i\right)A^{2m-n}}
\le \frac{1}{2m}\bigl(nA+(2m-n)A\bigr)=A.
\]
因此
\[
\left(\prod_{i=1}^n a_i\right)A^{2m-n}\le A^{2m}
\implies \prod_{i=1}^n a_i \le A^n,
\]
从而
\[
\sqrt[n]{\prod_{i=1}^n a_i}\le A=\frac{1}{n}\sum_{i=1}^n a_i.
\]

综上，命题得证. 
\end{proof}

\begin{theorem}[Young 不等式]
设 $p,q \in (1,+\infty)$ 且
\[
\frac{1}{p}+\frac{1}{q}=1.
\]
则对任意 $a,b>0$，有
\[
ab\le \frac{1}{p}a^p+\frac{1}{q}b^q.
\]
\end{theorem}

\begin{proof}
原图先将不等式化为
\[
1\le \frac{1}{p}a^{p-1}b^{-1}+\frac{1}{q}a^{-1}b^{q-1}.
\]

注意到
\[
\left(a^{p-1}b^{-1}\right)^{1/p}\left(a^{-1}b^{q-1}\right)^{1/q}
=a^{1-\frac{1}{p}-\frac{1}{q}}b^{-1+\frac{q-1}{q}}
=1.
\]

令
\[
A=\left(a^{p-1}b^{-1}\right)^{1/p},
\]
则
\[
a^{-1}b^{q-1}=\frac{1}{A^q}.
\]
所以只需证明
\[
1\le \frac{1}{p}A^p+\frac{1}{q}\frac{1}{A^q}.
\]

由对称性，不妨设 $A \ge 1$. 任取整数 $m>p$，令
\[
k=\left\lfloor \frac{m}{p}\right\rfloor+1.
\]
则
\[
k-1\le \frac{m}{p}<k
\implies
\frac{1}{p}\le \frac{k}{m}\le \frac{1}{p}+\frac{1}{m}.
\]
于是
\[
A^{m/k}\le A^p,\qquad \frac{1}{A^{m/(m-k)}}\le \frac{1}{A^q}.
\]
故
\[
\begin{aligned}
\frac{1}{p}A^p+\frac{1}{q}\frac{1}{A^q}
&\ge \frac{k}{m}A^{m/k}+\left(1-\frac{k}{m}\right)\frac{1}{A^{m/(m-k)}}-\frac{1}{m}A^p \\
&=\frac{1}{m}\left(kA^{m/k}+(m-k)\frac{1}{A^{m/(m-k)}}\right)-\frac{1}{m}A^p \\
&\ge \sqrt[m]{A^m\cdot A^{-m}}-\frac{1}{m}A^p \\
&=1-\frac{1}{m}A^p.
\end{aligned}
\]

令
\[
S=\frac{1}{p}A^p+\frac{1}{q}\frac{1}{A^q}-1.
\]
则对任意整数 $m>p$，都有
\[
S\ge -\frac{1}{m}A^p.
\]
令 $m \to \infty$，得 $S \ge 0$，从而
\[
1\le \frac{1}{p}A^p+\frac{1}{q}\frac{1}{A^q}.
\]
于是
\[
ab\le \frac{1}{p}a^p+\frac{1}{q}b^q.
\]
\end{proof}

\begin{theorem}
设 $x \ge 0$，则
\[
(1+x)^\alpha-1\le \alpha x,\qquad \forall \alpha \in (0,1],\ x \ge 0;
\]
\[
(1+x)^\alpha-1\ge \alpha x,\qquad \forall \alpha \ge 1,\ x \ge 0;
\]
\[
(1+x)^\alpha-1\ge \frac{\alpha x}{1+x},\qquad \forall \alpha \ge 0,\ x \ge 0.
\]
\end{theorem}
